\documentclass[11pt, letter]{article}
\usepackage[inner=1.5cm,outer=1.5cm,top=2.5cm,bottom=2.5cm]{geometry}
\pagestyle{empty}
\usepackage{graphicx}
\usepackage{fancyhdr, lastpage, bbding, pmboxdraw}
\usepackage[usenames,dvipsnames]{color}
\definecolor{darkblue}{rgb}{0,0,.6}
\definecolor{darkred}{rgb}{.7,0,0}
\definecolor{darkgreen}{rgb}{0,.6,0}
\definecolor{red}{rgb}{.98,0,0}
\usepackage[colorlinks,pagebackref,pdfusetitle,urlcolor=darkblue,citecolor=darkblue,linkcolor=darkred,bookmarksnumbered,plainpages=false]{hyperref}
\renewcommand{\thefootnote}{\fnsymbol{footnote}}
\pagestyle{fancyplain}
\fancyhf{}
\lhead{ \fancyplain{}{Astronomical Techniques} }
\rhead{ \fancyplain{}{\today} }
\fancyfoot[RO, LE] {page \thepage\ of \pageref{LastPage} }
\thispagestyle{plain}

\usepackage{listings}
\usepackage{caption}
\DeclareCaptionFont{white}{\color{white}}
\DeclareCaptionFormat{listing}{\colorbox{gray}{\parbox{\textwidth}{#1#2#3}}}
\captionsetup[lstlisting]{format=listing,labelfont=white,textfont=white}
\usepackage{verbatim}
\usepackage{fancyvrb}
\usepackage{acronym}
\usepackage{amsthm}
\VerbatimFootnotes

\definecolor{OliveGreen}{cmyk}{0.64,0,0.95,0.40}
\definecolor{CadetBlue}{cmyk}{0.62,0.57,0.23,0}
\definecolor{lightlightgray}{gray}{0.93}

% ---- heading helpers ----
% Major section rule (e.g. "Course Goals", "Grading Policy")
\newcommand{\sectionheading}[1]{%
  \vskip.25in
  \noindent\rule{\linewidth}{0.8pt}\\[2pt]
  \noindent{\large\textbf{#1}}\\[-6pt]
  \noindent\rule{\linewidth}{0.4pt}
  \vskip.1in}

% Goal sub-heading (numbered goal inside a section)
\newcommand{\goalheading}[2]{%
  \vskip.15in
  \noindent\textbf{\color{darkred}Goal #1: #2}
  \vskip.05in}

\lstset{
basicstyle=\ttfamily,
keywordstyle=\color{OliveGreen},
commentstyle=\color{gray},
numbers=left,
numberstyle=\tiny,
stepnumber=1,
numbersep=5pt,
backgroundcolor=\color{lightlightgray},
frame=none,
tabsize=2,
captionpos=t,
breaklines=true,
breakatwhitespace=false,
showspaces=false,
showtabs=false,
columns=flexible,
morekeywords={__global__, __device__},
}

%%%%%%%%%%%%%%%%%%%%%%%%%%%%%%%%%%%%
\begin{document}

% ============================================================
%  Title block
% ============================================================
\begin{center}
  {\Large \textsc{Astronomical Techniques}}\\[4pt]
  {\large ASTR 5110 \quad|\quad Fall 2027}
\end{center}

\vskip.2in
\noindent\rule{6in}{0.8pt}
\begin{center}
\begin{minipage}[t]{.75\textwidth}
\begin{tabular}{llcccll}
\textbf{Instructor:} & Yifan Zhou      & & & & \textbf{Time:}  & \\
\textbf{Email:}      & \href{mailto:yzhou@virginia.edu}{yzhou@virginia.edu}
                                        & & & & \textbf{Place:} & Astronomy Building \\
\textbf{Office Hours:} & & & & & &
\end{tabular}
\end{minipage}
\end{center}
\noindent\rule{6in}{0.8pt}

% ============================================================
\sectionheading{Course Description}
% ============================================================

This is a course in practical observational astronomy for graduate students in the Astronomy Department. The course includes substantial observing sessions and lab-based work. 

The course will consist of lectures, practical sessions, problem sets and observing projects (``labs'') with an emphasis on building practical observing experience. The main topics include understanding the night sky, designing observational experiments, planning and making observations, using and understanding typical modern astronomical instruments, reducing data, analyzing results, and reporting your results. The primary goal of the course is to train students in the basic working knowledge of the practical use of optical/near-infrared observing facilities, especially those that they may need for research on their thesis projects.

While this training is directly useful to students intending to undertake observationally oriented research projects and careers, the gained knowledge is also essential for theorists. Theorists must develop a basic understanding of how data are collected, reduced and analyzed to have a proper appreciation for the meaning, methods, limitations, and statistical interpretation of empirical work used to constrain and test theory. These skills are also needed to operate our department's telescope facilities—a skill you will likely need as a teaching assistant or participant in department public nights and outreach activities.

% [course description text here]

% ============================================================
\sectionheading{Course Goals and Learning Objectives}
% ============================================================

% ---- Goal 1 ------------------------------------------------
\goalheading{1}{Master the Signals in Your Data}

Cutting-edge astronomy demands not only knowing what your data say, but how reliably they say it.
By the end of this course, you will be able to:
\begin{itemize}
  \item Describe the fundamentals of telescope optical design, photon collection at optical and
        infrared wavelengths, and detector physics.
  \item Collect, reduce, and analyze astronomical datasets across three observing modes:
        imaging, spectroscopy, and interferometry.
  \item Quantify signals and uncertainties through hands-on data analysis activities.
\end{itemize}

% ---- Goal 2 ------------------------------------------------
\goalheading{2}{Formulate the Questions}

Good science begins with a well-posed question.
By the end of this course, you will be able to:
\begin{itemize}
  \item Connect current research questions in astronomy to the observational techniques used
        to address them.
  \item Select appropriate observing techniques for a given scientific question.
  \item Identify gaps in existing observational data that limit progress on specific astronomical
        problems.
\end{itemize}

% ---- Goal 3 ------------------------------------------------
\goalheading{3}{Design the Experiment}

Translating a question into a feasible observing program is a core professional skill.
By the end of this course, you will be able to:
\begin{itemize}
  \item Construct testable hypotheses from scientific questions using astronomical methods.
  \item Select observing methods, identify appropriate targets, and design an experiment.
  \item Justify your choices of observing mode, technique, and target list.
  \item Execute an observing program through hands-on laboratory work.
\end{itemize}

% ---- Goal 4 ------------------------------------------------
\goalheading{4}{Compete for Resources}

Astronomical facilities are oversubscribed.
Effective writing and presentation determine who gets access.
By the end of this course, you will be able to:
\begin{itemize}
  \item Write telescope proposals and scientific papers following professional standards.
  \item Provide constructive written peer review of scientific proposals.
  \item Present an observing program orally to a scientific audience.
\end{itemize}

% ---- Goal 5 ------------------------------------------------
\goalheading{5}{Apply These Skills Broadly}

The quantitative reasoning and experimental design habits you develop here transfer well beyond
astronomy.
By the end of this course, you will be able to:
\begin{itemize}
  \item Connect the hypothesis-testing and data analysis approaches from this course to
        problems in other scientific and technical domains.
\end{itemize}

% ============================================================
\sectionheading{Main References}
% ============================================================

This is a restricted list of books that will be touched during the course.
You should consult them occasionally.

\begin{itemize}
  \item % reference 1
\end{itemize}

\sectionheading{Night labs}

\sectionheading{Group work and projects}

\sectionheading{Observatory field trip}

% ============================================================
\sectionheading{Tentative Course Outline}
% ============================================================

\begin{center}
\begin{minipage}{5in}
\begin{flushleft}
{\color{darkgreen}{\Rectangle}}~
\end{flushleft}
\end{minipage}
\end{center}

% ============================================================
\sectionheading{Grading Policy}
% ============================================================

\begin{center}
\begin{tabular}{lc}
\hline
\textbf{Component}  & \textbf{Weight} \\
\hline
Homework \& Quizzes & 30\% \\
Midterm 1           & 20\% \\
Midterm 2           & 20\% \\
Final Exam          & 30\% \\
\hline
\end{tabular}
\end{center}

% ============================================================
\sectionheading{Important Dates}
% ============================================================

\begin{center}
\begin{minipage}{3.8in}
\begin{flushleft}
Midterm \#1  \dotfill   \\
Midterm \#2  \dotfill        \\
Final Exam   \dotfill             \\
\end{flushleft}
\end{minipage}
\end{center}

% ============================================================
\sectionheading{Course Policies}
% ============================================================

\end{document}